\section*{Overview}

SOS DP means ``sum over subsets'' DP, but the idea is broader than just sums. It is the standard transform when each
mask needs information from all of its submasks or all of its supermasks, and the naive \(O(4^n)\) approach is too
slow.

\section*{When to Use It}

Use SOS DP when:

\begin{itemize}
  \item the state is a bitmask,
  \item each mask wants to aggregate values over all submasks or all supermasks,
  \item \(n\) is small enough for \(2^n\) states but not small enough for nested subset enumeration on every mask.
\end{itemize}

\section*{Core Idea}

Instead of iterating through every \((mask, submask)\) pair independently, process the bits one by one. At each bit,
decide whether information should flow from the version without that bit to the version with that bit.

That transforms the brute-force subset aggregation into a dimension-by-dimension DP over the hypercube.

\section*{Key Insight}

The expensive part is repeated reuse of the same smaller masks. SOS DP avoids recomputing that reuse from scratch. It
folds one bit at a time, very much like a fast transform.

\section*{Operations / Main Technique}

\begin{itemize}
  \item initialize \(\texttt{f[mask]}\) with the base values,
  \item for each bit, propagate from masks missing that bit to masks containing it,
  \item read the transformed array as ``sum over all submasks'' or the corresponding variant.
\end{itemize}

\paragraph{Worked example.}
If \(\texttt{f[mask]}\) counts exact masks, the transformed array after SOS can store how many exact masks are subsets
of each query mask. That is the version used in many inclusion-style bitmask problems.

\section*{Correctness Intuition}

After processing the first \(b\) bits, each state already contains the aggregate over all choices of those bits that are
compatible with it. When the next bit is processed, the same invariant extends by one dimension. After all bits are
processed, every submask has been accounted for exactly once.

\section*{Complexity Analysis}

For \(n\) bits:
\begin{itemize}
  \item time: \(O(n 2^n)\),
  \item memory: \(O(2^n)\).
\end{itemize}

\section*{Implementation}

The code includes both:

\begin{itemize}
  \item subset-sum transform,
  \item superset-sum transform.
\end{itemize}

That is the pair I reach for most often.

\section*{Common Pitfalls}

\begin{itemize}
  \item Using SOS DP when ordinary submask enumeration is already fast enough.
  \item Mixing up the subset and superset directions.
  \item Forgetting that the state count still grows as \(2^n\), so the bit count must stay small.
\end{itemize}

\section*{Variants / Extensions}

\begin{itemize}
  \item Fast zeta transform and Mobius inversion on subsets.
  \item XOR / AND / OR convolution families.
  \item Bitset tricks when the values are boolean instead of numeric counts.
\end{itemize}

\section*{Practice Problems}

\begin{itemize}
  \item Count how many given masks are subsets of each query mask.
  \item Optimize over all supermasks of a state.
  \item Inclusion-style DP where many masks share the same submask contributions.
\end{itemize}

\section*{References}

\begin{itemize}
  \item \href{https://usaco.guide/plat/dp-sos?lang=py}{USACO Guide: Sum over Subsets DP}.
  \item \href{https://codeforces.com/catalog?locale=en}{Codeforces Catalog}.
  \item \href{https://oi-wiki.org/dp/sos/}{OI Wiki: SOS DP}.
\end{itemize}
